% --------------------------------- Preamble -----------------------------------
\documentclass[11pt, letterpaper, twocolumn]{article}

\usepackage[T1]{fontenc} % use 8-bit fonts
\usepackage[margin=0.75in, columnsep=0.2in, % 3.4 in columns
    headsep=0.28in, footskip=0.42in]{geometry} % paper sizes
\usepackage{amsmath, amsfonts, amssymb} % improved math
\usepackage{bm} % bold math (must be after amsmath)
\usepackage{upquote} % upright single quotes in verbatim
\usepackage{pgfplots} % includes tikz, xcolor, graphicx
\usepackage[colorlinks=true, allcolors=gray]{hyperref}

\raggedbottom % do not force stretching text to fill the page
\setlength{\emergencystretch}{\hsize} % avoid overfull lines
\overfullrule=1mm % black band to show overfull lines
\pdfsuppresswarningpagegroup=1 % remove pdf image warnings
\frenchspacing % consistent spacing between words
\newcommand{\ul}[1] % vectors
    {{}\mkern1mu\underline{\mkern-1mu#1\mkern-1mu}\mkern1mu}
\DeclareSymbolFont{euler}{U}{eur}{m}{n} % symbol font
\DeclareMathSymbol{\PI}{\mathord}{euler}{25} % up pi
\pgfplotsset{compat=newest} % use newest pgfplot version
\newcommand{\EE}[2]{\ensuremath{{{#1}\!\times\!10^{#2}}}}
% ------------------------------------------------------------------------------

\title{Sigma-Point Kalman Filtering for Drone State Estimation}
\author{}
\date{}

\begin{document}
\maketitle

\section*{Description}

This project implements a Sigma-Point Kalman Filter (SPKF), specfically a Square-Root 
Unscented Kalman Filter (UKF), for estimating the lateral motion of a drone. The system
dynamics are nonlinear and include coupled position, speed, roll, and heading states.
Noisy GPS position measurements are used to recursively estimate the full state
over time.

\subsection*{System Model}

The state vector is:
\[ \ul{x} = \begin{bmatrix}
    p_x \\
    p_y \\
    s \\
    \phi \\
    \psi
\end{bmatrix}
\]
where $(p_x,p_y)$ is the position of the drone, $s$ is the speed, $\phi$ is the roll
angle, and $\psi$ is the heading angle. 

The drone is modeled using coordinated-turn dynamics:
\[\begin{aligned}
\dot{\ul{x}} &= \begin{bmatrix}
            s*cos\psi & s*sin\psi & a & \omega & \frac{g}{s}sin\phi
        \end{bmatrix} ^T + \ul{\nu} \\
\ul{\nu} &\sim N(0,\bm{Q}) \\
\bm{Q} &= diag\begin{bmatrix}
            0.1^2 & 0.1^2 & 0.1^2 & (2\frac{\pi}{180})^2 & (2\frac{\pi}{180})^2
        \end{bmatrix}
\end{aligned}\]
where $a$ is the acceleration input, $\omega$ is the roll-rate input, and $g$
is acceleration of gravity $(9.8 m/s^2)$.

The drone has noisy GPS measurements of $(p_x, p_y)$ position:
\[\begin{aligned}
    \bm{z} &= \bm{H}\ul{x} + \ul{\eta} \\
\bm{H} &= \begin{bmatrix}
            1 & 0 & 0 & 0 & 0 \\
            0 & 1 & 0 & 0 & 0
        \end{bmatrix} \\
\ul{\eta} &\sim N(0,\bm{R}) \\
\bm{R} &= diag\begin{bmatrix}
            3^2 & 3^2
        \end{bmatrix}
\end{aligned}\]

\section*{Sigma-Point Kalman Filtering}

Sigma-Point Kalmaning Filters approximate the propagation of a Gaussian distribution through
nonlinear dynamics using a deterministic set of sample points. These sigma points are chosen
to capture the mean and covariance of the state distribution.

Each sigma point is propagated through the nonlinear system, and the transformed points are 
recombined to recover the mean and covariance. This approach avoids linearization and provides
improved accuracy compared to the Extended Kalman Filter.

\subsection*{Unscented Kalman Filter}

In this project, the sigma-point framework is implemented using the Unscented transform. The UKF
selects sigma points based on the current state estimate and covariance, propagates them through
the nonlinear dynamics, and computes weighted estimates of the predicted mean and covariance.

The accuracy of the UKF depends on selecting the optimal values for the parameters $\alpha$,
$\beta$, and $\kappa$. $\alpha$ controls the spread of sigma points around the mean, $\beta$
encodes prior knowledge about the distribution, and $\kappa$ is a secondary scaling parameter.

\subsection*{Square-Root Formulation}

A square-root form of the UKF is used in this implementation. Instead of propagating the covariance
matrix directly, the filter maintains and propagates its Cholesky factor.

This approach provides several advantages:
\begin{itemize}
    \item improved numerical stability
    \item preservation of positive definiteness
    \item reduced sensitivity to roundoff errors
\end{itemize}

These properties are especially important in nonlinear estimation problems, where covariance 
matrices can degrade over time when propagated directly.

\section*{Simulation Setup}

The filter is run with a sampling period of 0.1 s over 1000 samples. The initial state
estimate is
\[
\hat{\ul{x}} = \begin{bmatrix}
    0 & 0 & 25 & 0 & 0
\end{bmatrix}^{\top}
\]
with initial covariance
\[
\hat{\mathbf{P}} = \mathrm{diag}\left(\begin{bmatrix}
    10^2 & 10^2 & 3^2 & (\frac{1}{18})^2 & (\frac{1}{18})^2
\end{bmatrix}
    \right)
\]

\section*{Results}

Figure~\ref{fig:path} shows the true and estimated drone trajectory.

\begin{figure}[h]
    \centering
    \includegraphics[width=\linewidth]{../figures/fig_paths.pdf}
    \caption{True and estimated drone trajectory}
    \label{fig:path}
\end{figure}

\pagebreak

Figures~\ref{fig:perr}--\ref{fig:ryerr} show the position, speed, roll, and heading 
errors over time along with their corresponding $3-\sigma$ bounds.

\begin{figure}[h]
    \centering
    \includegraphics[width=\linewidth]{../figures/fig_position_errors.pdf}
    \caption{Position errors}
    \label{fig:perr}
\end{figure}

\begin{figure}[h]
    \centering
    \includegraphics[width=\linewidth]{../figures/fig_speed_error.pdf}
    \caption{Speed errors}
    \label{fig:serr}
\end{figure}

\begin{figure}[h]
    \centering
    \includegraphics[width=\linewidth]{../figures/fig_angle_errors.pdf}
    \caption{Roll and Heading errors}
    \label{fig:ryerr}
\end{figure}

\pagebreak

\section*{Discussion}

The results demonstrate the sigma-point filtering is effective for nonlinear flight dynamics.
Even with limited measurements, the filter is able to infer the full state through the
structure of the system dynamics.

The square-root formulation improves numerical robustness and helps maintain a consistent
covariance representation over time. This is particularly important for nonlinear systems
and longer simulations.

This project highlights the advantage of sigma-point methods: accurate nonlinear estimation
without requiring analytical linearization, combined with strong numerical properties when
implemented in square-root form.

\end{document}